\documentclass[10pt]{article}
\usepackage[a4paper,margin=2.5cm]{geometry}
\usepackage{graphicx}
\usepackage{amsmath}
\usepackage[font=small,labelfont=bf]{caption}
\graphicspath{{figures/}}

\title{A synthetic manuscript for the braindance examples}
\author{}
\date{}

\begin{document}
\maketitle

All data in this document are synthetic, made by \texttt{make\_data.py}.
Figure~\ref{fig:overview} shows the study, Fig.~\ref{fig:convergence} the grid convergence,
Fig.~\ref{fig:results} the results per case and Fig.~\ref{fig:sweep} a hyperparameter sweep.

\begin{figure}[htbp]
  \centering
  \includegraphics[width=\linewidth]{overview}
  \caption{Overview of the study. (a) Instantaneous temperature of the two validation cases.
  (b) Training cases on a $3\times3$ grid of Reynolds number $Re$ and heating ratio $q$, and the two validation cases.
  (c) Mean temperature profiles of all cases.}
  \label{fig:overview}
\end{figure}

\begin{figure}[htbp]
  \centering
  \includegraphics[width=0.8\linewidth]{convergence}
  \caption{Grid convergence of three schemes: $L_2$ error against the number of grid points $N$,
  mean and range over five random initial conditions.}
  \label{fig:convergence}
\end{figure}

\begin{figure}[htbp]
  \centering
  \includegraphics[width=0.7\linewidth]{results_table}
  \caption{Error of the baseline and the proposed method for each case, the gain
  $\Delta=\varepsilon_\mathrm{baseline}-\varepsilon_\mathrm{proposed}$ and the run time.}
  \label{fig:results}
\end{figure}

\begin{figure}[htbp]
  \centering
  \includegraphics[width=\linewidth]{sweep}
  \caption{Gain over the baseline, in percentage points, of two models across learning rate and network width.}
  \label{fig:sweep}
\end{figure}

\end{document}
